\documentclass[10pt,a4paper]{amsart}
\usepackage[margin=19mm]{geometry}
\usepackage{amsmath,amssymb,mathtools}
\usepackage[scr=rsfs,cal=euler]{mathalfa}
\usepackage{xfrac}
\usepackage[unicode,hidelinks,bookmarks=false]{hyperref}
\newcommand{\ie}{i.e.\ }
\newcommand{\cf}{cf.\ }
\newcommand{\resp}{respectively\ }
\newcommand{\Subsection}{\subsection}
\providecommand{\nobreakdash}{\nobreak\hbox{-}}
\setlength{\emergencystretch}{2em}
\multlinegap0pt
\usepackage{bm}
\usepackage{bbm}
\usepackage{dsfont}
\usepackage{xspace}
\usepackage{cancel}
\usepackage{mathtools}
\usepackage{amsthm}
\makeatletter
\@ifundefined{thm}{\newtheorem{thm}{Theorem}}{}
\@ifundefined{defn}{\newtheorem{defn}[thm]{Definition}}{}
\@ifundefined{lem}{\newtheorem{lem}[thm]{Lemma}}{}
\@ifundefined{prop}{\newtheorem{prop}[thm]{Proposition}}{}
\makeatother
\newcommand{\RR}{\mathbb{R}}
\newcommand{\eps}{\varepsilon}
\newcommand{\itemizeEqnVSpacing}{\rule{0pt}{1pt}\vspace*{-12pt}}
\title[On Maximal Functions Associated to Families of Curves in the]{On Maximal Functions Associated to Families of Curves in the Plane}
\author{Joshua Zahl}
\date{Source: arXiv:2307.05894v3 (CC BY 4.0)}
\begin{document}
\maketitle

\noindent\emph{Source and licence.} On Maximal Functions Associated to Families of Curves in the Plane. Version arXiv:2307.05894v3 (CC BY 4.0), distributed under the Creative Commons Attribution 4.0 International licence, as recorded in the arXiv licence metadata for this exact version and shown on the abstract page https://arxiv.org/abs/2307.05894v3. Full rights notice: licenses/arxiv-2307.05894-CC-BY-4.0.txt.\par\smallskip

\noindent\textit{Editorial scope.} Counted unit: the Prop whose source label is \texttt{maximalFnBdNonconcentrated} in the article (source0.tex lines 2152--2165, source label \texttt{maximalFnBdNonconcentrated}), delivered together with the article's own complete original proof of it (source0.tex lines 2167--2301, one \texttt{proof} environment of 135 lines). No mathematics is written, repaired or re-derived: statement and proof are the article's own text line for line. Required context retained uncounted: RectBoundThm (lines 560--567); defnTangencyPrism (lines 702--711); tangencyPrismVsRect (lines 728--730); pullbackRect (lines 759--785); CkNormRectangleRescaling (lines 789--791); rectangleDecompProp (lines 1821--1858); inductionPropHypothesis (lines 1827--1829); conclusionB (lines 1848--1850); inductionPropHypothesisAfterIteration (lines 1854--1856). Every definition, display and label of those lines is retained unchanged. Omitted from this package and not redistributed: 1--559, 568--701, 712--727, 731--758, 786--788, 792--1820, 1830--1847, 1851--1853, 1857--2151, 2166--2166, 2302--2961. Nothing inside a retained excerpt is omitted or altered: every retained body is delivered exactly as the article's lines are written, so no omission receipt of any kind accompanies this package. Citations: the retained excerpts carry no citation command at all, so no bibliography entry of the article is transcribed and no reference is redistributed. Reference adaptation: none was needed and none was made; every reference inside the delivered unit resolves inside it, so no unresolved reference or citation is produced. Printed slips kept literal: no printed word, symbol, hypothesis or number of the counted statement or of its proof was altered. Layout and class: the article's own class is \texttt{article} with packages including geometry, color, latexsym, amsmath, amssymb, amsthm, graphicx, subfigure, revsymb, enumerate, xcolor, mathtools, nth; the wrapper is the lane's pinned \texttt{amsart} editorial wrapper at 10pt on A4, which restates the theorem-like environments the retained text needs and adds the article's own macro definitions that the retained text needs (\texttt{\textbackslash RR}, \texttt{\textbackslash eps}, \texttt{\textbackslash itemizeEqnVSpacing}), with extra wrapper packages (bm, bbm, dsfont, xspace, cancel, mathtools). The delivered numbering is the unit's own. Assets: no retained span contains a figure environment, a graphics-inclusion command, a PDF-page inclusion, a table or a TikZ picture; no image file is copied into this package, the package declares no asset dependency and no image file is redistributed with it. Licence: version arXiv:2307.05894v3 of this article is distributed under the Creative Commons Attribution 4.0 International licence, as recorded in the arXiv licence metadata for this exact version and shown on the abstract page https://arxiv.org/abs/2307.05894v3. A full rights notice is delivered at licenses/arxiv-2307.05894-CC-BY-4.0.txt. Characters: every retained excerpt is pure ASCII, so the delivered engine, XeLaTeX with its default Latin Modern text font, reports no missing character.
\begin{thm}\label{RectBoundThm}
Let $k,N\geq 1$ and $\eps>0$. Then there exists $\eta,\delta_0>0$ so that the following holds for all $\delta\in(0,\delta_0]$. 

Let $F$ be a set of (univariate) polynomials of degree at most $\delta^{-\eta}$, each of which has $C^{k}$ norm at most 1. Suppose that $\#F\leq\delta^{-N}$. Let $\mu\geq 1$ and let $\mathcal{R}$ be a set of pairwise incomparable $(\delta,k)$ rectangles that are $\mu$-rich and $\eps$-robustly broad with error at most $\delta^{-\eta}$ with respect to $F$. Then
\begin{equation}\label{RectBdEqn}
\# \mathcal{R} \leq \delta^{-\eps} \Big(\frac{\# F}{\mu}\Big)^{\frac{k+1}{k}}.
\end{equation}
\end{thm}

\begin{defn}\label{defnTangencyPrism}
A $(\delta;k)$ \emph{tangency prism} is a set $P$ of the form
\begin{equation}\label{tangencyPrismEqn}
\begin{split}
\Big\{(t, y_0,\ldots&,y_{k-1})\in\RR^{k+1}\colon  t\in [a, a+\delta^{1/k}],\\
&\Big| y_j - \sum_{i=j}^{k-1}\frac{(t-a)^{i-j}}{(i-j)!}b_i\Big|\leq K\delta^{1-j/k},\ j=0,\ldots,k-1\Big\}.
\end{split}
\end{equation}
In the above expression, $a\in[0,1]$ and $b_0,\ldots,b_{k-1}\in [-1,1]$ are parameters that define the tangency prism, and $K$ is a constant depending on $k$; the specific choice of $K$ will be fixed in Lemma \ref{tangencyPrismVsRect} below. We call $I = [a, a+\delta^{1/k}]$ the \emph{interval associated to} $P$.
\end{defn}

\begin{lem}\label{tangencyPrismVsRect}
If the quantity $K=O(1)$ from Definition \ref{defnTangencyPrism} is chosen appropriately (depending on $k$), then the following is true. Let $R$ be a $(\delta,k)$ tangency rectangle, let $f$ be a function with $C^k$ norm at most 1, and suppose that $f\sim R$. Then $\mathcal{J}_{k-1}f\sim \hat R$.
\end{lem}

\begin{defn}\label{pullbackRect}
Let $\rho>0$ and let $S=S^g(I)$ be a $(\rho;k)$ rectangle; here $\Vert g \Vert_{C^k}\leq 1$ and $I=[a, a+\rho^{1/k}]$. Let $K=K(k)\geq 1$ be the constant from Definition \ref{defnTangencyPrism}, and let $c=\big((k+1)K\big)^{-1}$. 

\smallskip

\noindent (A): For $x\in I$, define
\[
\phi^S(x,y) = \big( \rho^{-1/k}(x-a),\ c\rho^{-1}(y-g(x))\big).
\]
With this definition, we have $\phi^S(S)=[0,1]\times [-c,c]$.

\smallskip

\noindent (B): For $f\colon [0,1]\to\RR$, define $f_S$ to be the function whose graph is $\phi^S(\operatorname{graph}f|_I)$. 

\smallskip

\noindent (C): For $x\in I$ and $y_0,\ldots,y_{k-1}\in\RR^k$, define 
\begin{equation*}
\begin{split}
\psi^S(x,&y_0,\ldots,y_{k-1}) \\
= \Big(&\rho^{-1/k}(x-a),\ c\rho^{-1}(y_0-g(x)),\ c\rho^{-1+1/k}(y_1-g'(x)),\\
& \qquad c\rho^{-1+2/k}(y_2-g''(x)),\ldots, c\rho^{-1/k}(y_{k-1}-g^{(k-1)}(x))\Big).
\end{split}
\end{equation*}
With this definition, we have $\psi^S(\hat S)=[0,1]\times [-Kc,Kc]^k = [0,1]\times [\frac{-1}{k+1},\frac{1}{k+1}]^k$.
\end{defn}

\begin{lem}\label{CkNormRectangleRescaling}
Let $S$ be a $(\rho;k)$ rectangle, let $\Vert f\Vert_{C^k}\leq 1$, and suppose $f\sim S$. Then $\Vert f_S \Vert_{C^k}\leq 1$. 
\end{lem}

\begin{prop}\label{rectangleDecompProp}
Let $k,N\geq 2$ and let $0<\eps,\eta\leq 1$. Then there exists $\delta_0>0$ such that the following is true for all $\delta\in (0,\delta_0]$. Let $F$ be a set of functions, each of which has $C^k$ norm at most $1$. Suppose furthermore that $\#F \leq \delta^{-N}$. For each $f\in F$, let $Y(f)\subset f^\delta$ be a shading of $f$. 


Suppose that the functions in $F$ are broad with respect to $(\delta,k)$ tangency rectangles, in the following sense: for every $x\in[0,1]\times[-1,1]$, every $\rho\geq\delta$, every $T\in [1, 1/\rho],$ and every $(\rho; k; T)$ rectangle $R$ containing $x$, we have
%
\begin{equation}\label{inductionPropHypothesis}
\#\{f\in F\colon x\in Y(f),\ f\sim R\}\leq\delta^{-\eta}T^{-\eps}\#\{f\in F\colon x\in Y(f)\}. 
\end{equation}

Then there exists the following:
\begin{itemize}
	\item A sub-shading $Y'(f)\subset Y(f)$ for each $f\in F$.
	\item A scale $\delta'\in [\delta,1]$.
	\item A set $\mathcal{R}$ of pairwise incomparable $(\delta'; k)$ rectangles.
	\item A number $\mu$, and for each $R\in\mathcal{R}$, a set $F(R)\subset \{f\in F\colon f\sim R\}$ of size $\mu$.
\end{itemize}
These objects have the following properties:
\begin{itemize}
\item[(A)] \itemizeEqnVSpacing
\begin{equation}\label{LkBdControlledBySumOverRect}
\int_{[0,1]\times[-1,1]}\Big(\sum_{f\in F}\chi_{Y(f)}\Big)^{\frac{k+1}{k}}\leq \delta^{-O(\sqrt\eta)}\sum_{R\in\mathcal{R}} \int_{R}\Big(\sum_{f\in F(R)}\chi_{Y'(f)}\Big)^{\frac{k+1}{k}}.
\end{equation}

%
\item[(B)]
Either $\#\mathcal{R}=1$, or for every $R\in\mathcal{R}$, every $\rho\in[\delta',1]$, every $T\in [1,1/\rho]$, and every $(\rho; k; T)$ rectangle $R'\supset R$, we have
\begin{equation}\label{conclusionB}
\#\{f\in F(R)\colon f\sim R'\} \leq (\delta')^{-2\sqrt\eta}T^{-\eps} \mu.
\end{equation}
%
\item[(C)]
For every $R\in\mathcal{R}$, let $\tilde F(R)=\{f_R\colon f\in F(R)\}$, and for every $\tilde f\in \tilde F(R)$, let $\tilde Y'(\tilde f)=\phi^R(Y'(f)\cap R)$ (recall Definition \ref{pullbackRect}). Then for every point $x\in[0,1]\times[-1,1]$, every $\rho\geq\delta/\delta'$, every $T\in[1,1/\rho],$ and every $(\rho; k-1; T)$ rectangle $R^\dag$ containing $x$, we have
\begin{equation}\label{inductionPropHypothesisAfterIteration}
\#\{\tilde f\in \tilde F(R)\colon x\in \tilde Y'(\tilde f),\  \tilde f\sim R^\dag\}\lessapprox_\eps  T^{-\eta/2} \#\{\tilde f\in \tilde F(R)\colon x\in \tilde Y'(\tilde f)\}. 
\end{equation}
\end{itemize}
\end{prop}

\begin{prop}\label{maximalFnBdNonconcentrated}
Let $k,N\geq 1,$ $0\leq \alpha\leq 1$, and let $\eps>0$. Then there exists $\eta>0$ such that the following is true for all $\delta,\tau>0$. Let $F$ be a set of polynomials of degree at most $\delta^{-\eta}$, each of which has $C^{k}$ norm at most 1. Suppose that $\#F\leq\delta^{-N}$. For each $f\in F$, let $Y(f)$ be a $\delta$--thick shading striped by a $(\tau,\alpha;\delta^{-\eta})$-set. Suppose that for all $x\in [0,1]^2$, all $\rho\in[\delta,1]$, all $T\in [1, 1/\rho]$, and all $(\rho; k; T)$ rectangles $R$ containing $x$, we have
%
\begin{equation}\label{rectangleConcentrationBdNonconcentratedProp}
\#\{f\in F\colon x\in Y(f),\ f\sim R\} \leq \delta^{-\eta}T^{-\eps}\#\{f\in F\colon x\in Y(f)\}.
\end{equation}
%

Then
\begin{equation}\label{NonconcentratedPropLpBound}
\Big \Vert \sum_{f\in F}\chi_{Y(f)}\Big\Vert_{\frac{k+1}{k}}\lesssim_\eps  \delta^{-\eps}\delta^{\frac{k+\alpha}{k+1}}\tau^{\frac{k(1-\alpha)}{k+1}}
(\#F).
\end{equation}
\end{prop}

\begin{proof}
We prove the result by induction on $k$. 

\medskip

\noindent{\textbf{The base case}.}\ \ We begin with the base case $k=1$. By dyadic pigeonholing we can find a number $\mu$ and a set $X\subset[0,1]^2$ so that
\[
\Big \Vert \sum_{f\in F}\chi_{Y(f)}\Big\Vert_2^2\lesssim (\log 1/\delta) \mu^2|X|, \quad\textrm{and}\quad \mu\leq \sum_{f\in F}\chi_{Y(f)}(x)< 2\mu\quad\textrm{for all}\ x\in X.
\]
%
Thus in order to establish \eqref{NonconcentratedPropLpBound}, it suffices to show that
\begin{equation}\label{boundOnSizeOfX}
|X| \lesssim_\eps \delta^{-\eps/2}\delta^{1+\alpha}\tau^{1-\alpha}\mu^{-2}(\#F)^2.
\end{equation}


By \eqref{rectangleConcentrationBdNonconcentratedProp} with $\rho=2\delta$ and $T=2^{1/\eps}\delta^{-\eta/\eps}$, we have that for each $x\in X$, there are $\gtrsim\mu^2$ pairs $f,g\in F$ with the following two properties:
\begin{itemize}
\item[(i)] $x\in Y(f)\cap Y(g)$.
\item[(ii)] The connected component of $f^\delta\cap g^\delta$ containing $x$ projects to an interval of length at most $2^{1/\eps}\delta^{1-\eta/\eps}$ on the $x_1$-axis.
\end{itemize} 
For each pair $f,g\in F$, let 
\begin{equation*}
\begin{split}
W(f,g) = \{x\in & f^\delta\cap g^\delta\colon  \textrm{the connected component of}\ f^\delta\cap g^\delta\ \textrm{containing}\ x \\
& \textrm{projects to an interval of length}\ \leq 2^{1/\eps}\delta^{1-\eta/\eps}\ \textrm{on the $x_1$-axis}\}.
\end{split}
\end{equation*}
Since $f$ and $g$ are polynomials of degree at most $\delta^{-\eta}$, $f^\delta\cap g^\delta$ is a union of $O(\delta^{-\eta})$ connected components, and thus the projection of $W(f,g)$ to the $x_1$-axis is a union of $O(\delta^{-\eta})$ intervals, each of length at most $2^{1/\eps}\delta^{1-\eta/\eps-\eta}\lesssim_\eps \delta^{1-2\eta/\eps}$. Denote this set by $I(f,g)\subset[0,1]$.

Since $Y(f)$ is a $\delta$--thick shading striped by a $(\tau,\alpha;\delta^{-\eta})$-set, we have
\begin{equation}\label{Y1FcapIfg}
\begin{split}
|Y(f) \cap (I(f,g)\times\RR)|& \leq \sum_{I\subset I(f,g)} |Y(f) \cap (I\times\RR)|\\
&\lesssim_\eps \delta^{-\eta} \left\{ \begin{array}{ll} 
(\delta \tau) \big(\delta^{-\eta} (\delta^{1-2\eta/\eps}/\tau)^\alpha\big),& \delta^{1-2\eta/\eps}\geq\tau\\
\delta^{2-2\eta/\eps},& \delta^{1-2\eta/\eps}\leq\tau
\end{array}
\right.\\
&\leq \delta^{1+\alpha-3\eta/\eps}\tau^{1-\alpha},
\end{split}
\end{equation}
where in the first inequality, the sum is taken over the connected components (i.e.~intervals) in $I(f,g)$. This sum has $O(\delta^{-\eta})$ terms.




Let $\mathcal{T}$ be the set of triples $(x, f,g)\in X\times F^2$, where the pair $f,g$ satisfies Items (i) and (ii). Then $|\mathcal{T}|\sim \mu^2 X$, where $|\cdot|$ denotes the product of two-dimensional Lebesgue measure on $X$ and counting measure on $F^2$.

For each $f,g\in F$, we have
\[
|\{x\in X\colon (x, f,g)\in \mathcal{T}\}| \leq  |Y(f) \cap (I(f,g)\times\RR)|\leq \delta^{1+\alpha-3\eta/\eps}\tau^{1-\alpha},
\]
where the final inequality used \eqref{Y1FcapIfg}. We conclude that 
%
\[
|\mathcal{T}|\lesssim_\eps  \delta^{1+\alpha-3\eta/\eps}\tau^{1-\alpha}(\#F)^2.
\]
On the other hand, we have $|X|\leq \mu^{-2}|\mathcal{T}|$, and thus
\begin{equation}\label{boundOnX}
|X|\lesssim_\eps  \delta^{1+\alpha-3\eta/\eps}\tau^{1-\alpha}\Big(\frac{\#F}{\mu}\Big)^2.
\end{equation}
If we select $\eta\leq \eps^2/6$, then \eqref{boundOnX} implies \eqref{boundOnSizeOfX}. This completes our proof of the base case.

\medskip

\noindent \textbf{The induction step.}\ 
Suppose that $k\geq 2$ and that the result has been established for $k-1$. Fix $0\leq\alpha\leq 1$ and $\eps>0$. Let $\eta>0$ be a quantity to be specified below, let $\delta,\tau>0$, and let $F$ and $Y(f)$ satisfy the hypotheses of Proposition \ref{maximalFnBdNonconcentrated} with this value of $\eta$. First, let $\delta_0>0$ be a small quantity to be chosen below, which depends on $k$ and $\eps$. We may suppose that $\delta\leq\delta_0$, since otherwise \eqref{NonconcentratedPropLpBound} is trivial, provided we choose the implicit constant sufficiently large.  

Let $\eta_1=\eta_1(k,\eps)$ be the output from Theorem \ref{RectBoundThm}, with $\eps/2$ in place of $\eps$. Let $\eta_2 = \eta_1^2/4$. We will select $\eta>0$ sufficiently small so that $\eta\leq\eta_2$. Thus the shadings $\{Y(f)\colon f\in F\}$ satisfy Hypothesis \eqref{inductionPropHypothesis} of Proposition \ref{rectangleDecompProp} with $\eps/2$ in place of $\eps$ (this is because decreasing $\eps$ can only make \eqref{inductionPropHypothesis} easier to satisfy) and $\eta_2$ in place of $\eta$. Applying Proposition \ref{rectangleDecompProp}, we get a sub-shading $Y'(f)\subset Y(f)$; a scale $\delta'\in [\delta,1]$; a set $\mathcal{R}$ of $(\delta', k)$ rectangles; sets $F(R),$ $R\in\mathcal{R}$; and a multiplicity $\mu\leq\#F$. 


By Proposition \ref{rectangleDecompProp} Item (B), either $\#\mathcal{R}=1$, or we can apply Theorem \ref{RectBoundThm} (recall that we selected $\eta_1$ sufficiently small to ensure that Theorem \ref{RectBoundThm} can be applied, and $\eta_2 = \eta_1^2/4$; c.f.~\eqref{conclusionB}) to conclude that
\begin{equation}\label{rectangleBdInPropMaximalFnBdNonconcentrated}
\# \mathcal{R} \leq \delta^{-\eps/2}  \Big(\frac{\# F}{\mu}\Big)^{\frac{k+1}{k}}.
\end{equation}
Technically, we can only apply Theorem \ref{RectBoundThm} if $\delta$ is sufficiently small in terms of $k$ and $\eps$. However if this is not the case then \eqref{rectangleBdInPropMaximalFnBdNonconcentrated} follows from the crude bound $\# \mathcal{R}\lesssim (\delta')^{-k-1}\lesssim_{\eps}1$.


We next explore the consequences of Item (C) from Proposition \ref{rectangleDecompProp}. We first consider the case where $\delta'\leq\delta^{1-\eps/2}$. By Item (A) from Proposition \ref{rectangleDecompProp}, we have 
\begin{equation}
\begin{split}
\Big \Vert \sum_{f\in F}\chi_{Y(f)}\Big\Vert_{\frac{k+1}{k}}^{\frac{k+1}{k}} 
& \leq \delta^{-O(\sqrt \eta_2)}\sum_{R\in\mathcal{R}} \int_{R}\Big(\sum_{f\in F(R)}\chi_{Y'(f)}\Big)^{\frac{k+1}{k}}\\
& \lesssim \left\{ \begin{array}{ll} \delta^{-O(\eta_1)}(\#\mathcal{R}) \mu^{\frac{k+1}{k}} (\delta\tau)\Big( \delta^{-\eta}(\delta^{1/k}/\tau)^{\alpha}\Big), & \tau\leq\delta^{1/k} \\ \delta^{-O(\eta_1)}(\#\mathcal{R}) \mu^{\frac{k+1}{k}} \delta^{\frac{k+1}{k}}, & \tau>\delta^{1/k} \end{array}\right. \\
& \lesssim\delta^{-\eps/2 -O(\eta_1)}\delta^{\frac{k+\alpha}{k}}\tau^{1-\alpha}  (\#F)^{\frac{k+1}{k}},
\end{split}
\end{equation}
and we have established \eqref{NonconcentratedPropLpBound} and completed the proof, provided we select $\eta_1$ sufficiently small depending on $\eps$ and $k$, and provided we select the implicit constant in \eqref{NonconcentratedPropLpBound} sufficiently large depending on $\eps$ and $k$.

For the remainder of the proof we consider the case where $\delta'>\delta^{1-\eps/2},$ so in particular the implicit constant $O_{\eps}(\log 1/\delta)^{O_\eps(1)}$ from \eqref{inductionPropHypothesisAfterIteration} is bounded by $O_{\eps}(\log \frac{1}{\delta/\delta'})^{O_\eps(1)}$. For each $R\in\mathcal{R}$, let $\tilde F(R)$ and $\tilde Y'(f)$ be as defined in Item (C) of Proposition \ref{rectangleDecompProp}; these sets satisfy \eqref{inductionPropHypothesisAfterIteration}.

If $\delta_0$ (and thus $\delta/\delta'$) is sufficiently small, then $\tilde F(R)$ and $\tilde Y'(f)$ will satisfy the induction hypothesis \eqref{rectangleConcentrationBdNonconcentratedProp} with the parameters changed as follows:
\begin{itemize}
\item $k$ is replaced by $k-1$.
\item $\delta$ is replaced by $\tilde\delta = \delta/\delta'$.
\item $\tau$ is replaced by $\tilde\tau = \tau/(\delta')^{1/k}$
\item The functions $f\in F$ are polynomials of degree $\delta^{-\eta} \leq \tilde\delta^{-2\eta/\eps}$, each of which have $C^k$ norm at most 1 (the latter property follows from Lemma \ref{CkNormRectangleRescaling}).
\item The shadings $\tilde Y'(f)$ are $\tilde\delta$--thick shadings striped by a $(\tilde\tau,\alpha;\tilde\delta^{-2\eta/\eps})$-set.
\item The shadings $\tilde Y'(f)$ satisfy \eqref{rectangleConcentrationBdNonconcentratedProp}, with $T^{-\eta_2/2}$ in place of $T^{-\eps}$, and $O_{\eps}(\log 1/\tilde\delta)^{O_\eps(1)}$ in place of $\delta^{-\eta}$.
\end{itemize}

It is now time to apply the induction hypothesis: we apply Proposition \ref{maximalFnBdNonconcentrated} with $k-1$ in place of $k$; $\alpha$ unchanged; and $\eta_2/2$ in place of $\eps$. Let $\eta_3$ be the output from this proposition. If $\eta>0$ is selected sufficiently small, then $2\eta/\eps\leq\eta_3$. This means that the functions $f\in F$ are polynomials of degree at most $\tilde\delta^{-\eta_3}$, and the shadings $\tilde Y'(f)$ are striped by $(\tilde\tau,\alpha;\tilde\delta^{-\eta_3})$-sets. If $\delta_0$ and thus $\tilde\delta$ are sufficiently small, then the quantity $O_{\eps}(\log 1/\tilde\delta)^{O_\eps(1)}$ from the final item above is at most $\tilde\delta^{-\eta_3}$. Thus we can use the induction hypotheses to conclude that 
\begin{equation}\label{Lkm1Bd}
\Big \Vert \sum_{\tilde f\in \tilde F(R) }\chi_{\tilde Y'(\tilde f)}\Big\Vert_{\frac{k}{k-1}}\lesssim \delta^{-\eta_2/2}\tilde\delta^{\frac{k-1+\alpha}{k}} \tilde\tau^{\frac{(1-\alpha)(k-1)}{k}} (\#F(R)).
\end{equation}
We also have the $L^1$ bound
\begin{equation}\label{L1Bd}
\Big \Vert \sum_{\tilde f\in \tilde F(R) }\chi_{\tilde Y'(\tilde f)}\Big\Vert_{1} \leq \delta^{-\eta}\tilde\delta\tilde\tau^{1-\alpha}(\#F(R)).
\end{equation}
Interpolating \eqref{Lkm1Bd} and \eqref{L1Bd} using H\"older and recalling the definition of $\tilde\delta$ and $\tilde\tau$ (and the fact that $\eta\leq\eta_2$), we have
\begin{equation*}
\begin{split}
\Big \Vert \sum_{\tilde f\in \tilde F(R) }\chi_{\tilde Y'(\tilde f)}\Big\Vert_{\frac{k+1}{k}}^{\frac{k+1}{k}} 
& \leq \Big \Vert \sum_{\tilde f\in \tilde F(R) }\chi_{\tilde Y'(\tilde f)}\Big\Vert_{1}^{\frac{1}{k}} \Big \Vert \sum_{\tilde f\in \tilde F(R) }\chi_{\tilde Y'(\tilde f)}\Big\Vert_{\frac{k}{k-1}}\\
& \leq \delta^{-\eta_2}(\delta')^{-\frac{k+1}{k}} \delta^{\frac{k+\alpha}{k}}\tau^{1-\alpha}(\#F).
\end{split}
\end{equation*}

Undoing the scaling that maps $R$ to the unit square (this scaling distorts volumes by a factor of $(\delta')^{\frac{k+1}{k}}$), and recalling that $\#F(R)=\mu$ for each $R\in\mathcal{R}$, we conclude that
\begin{equation}\label{inductionHypothesesRectBd}
\int_R \Big( \sum_{f\in  F(R) }\chi_{Y'(f)}  \Big)^{\frac{k+1}{k}} \lesssim 
 \delta^{-\eta_2}\delta^{\frac{k+\alpha}{k}}\tau^{1-\alpha} \mu^{\frac{k+1}{k}}.
\end{equation}
Combining Item (A) from Proposition \ref{rectangleDecompProp} and \eqref{inductionHypothesesRectBd}, we conclude that 
%
\begin{equation}\label{controlOFLkNormInductionClose}
\begin{split}
\Big \Vert \sum_{f\in F}\chi_{Y(f)}\Big\Vert_{\frac{k+1}{k}}^{\frac{k+1}{k}} & \lesssim_\eps \delta^{-O(\sqrt\eta_2)}\sum_{R\in\mathcal{R}} \int_{R}\Big(\sum_{f\in F(R)}\chi_{Y'(f)}\Big)^{\frac{k+1}{k}}\\
& \lesssim  \delta^{-O(\eta_1)-\eta_2}(\#\mathcal{R})\delta^{\frac{k+\alpha}{k}}\tau^{1-\alpha} \mu^{\frac{k+1}{k}}\\
& \leq \delta^{-O(\eta_1)-\eta_2-\eps/2}\delta^{\frac{k+\alpha}{k}}\tau^{1-\alpha} (\# F)^{\frac{k+1}{k}}.
\end{split}
\end{equation} 
If we select $\eta_1$ and $\eta_2$ sufficiently small (depending on $\eps$ and $k$), then the term $\delta^{-O(\eta_1)-\eta_2-\eps/2}$ has size at most $\delta^{-\eps}$. This establishes \eqref{NonconcentratedPropLpBound} and closes the induction. 
\end{proof}

\end{document}
